\section{Structural Hypotheses Beyond the Replayed Kernel}\label{app:horizon}
\status{finite cardinal identities and explicitly conditional construction targets}
\subsection{The doubled Klein register and a possible reference fiber}
The group identity $V_4\times V_4\cong\F_2^4$ gives sixteen vectors. Over $\F_2$ each nonzero vector represents a unique projective point, so
\[
\F_2^4=\{0\}\sqcup\bigl(\F_2^4\setminus\{0\}\bigr),
\qquad |\mathrm{PG}(3,2)|=15.
\]
Consequently the \emph{sets} $V(G_{60})\times\F_2^4$ and $V(G_{60})\times(\F_2^4\setminus\{0\})$ have sizes $960$ and $900$. The omitted zero fiber has size $60$. This is the proposed ``objective handle'' bookkeeping: a complete sixteen-position register includes a distinguished reference fiber, while a punctured register retains fifteen positions.

No adjacency, current law or gluing operation follows from these counts. In particular $960=|\Aut(\mathcal K_{120})|$ is a group order, not an identification of that group with the proposed $960$-state set. An equivariant incidence map from the actual fifteen-position $L(P)$ carrier to the nonzero binary vectors must be supplied before identifying the two carriers. A further edge-level construction must reproduce the intended $G_{900}$ wiring and its transport labels. Neither claim is made by the present replay.

\subsection{A selected binary line is not yet physical time}
The finite vector space has dimension four and its projective space has dimension three. These are dimensions over a finite field, not a signature $(3,1)$ metric. A specified nonzero vector $t$ determines a line $T=\langle t\rangle$ and a canonical three-dimensional quotient $\F_2^4/T$. A direct-sum choice
\[
\F_2^4=D\oplus T,\qquad \dim D=3,
\]
requires a complement: there are eight complementary hyperplanes to a fixed $T$. Thus choosing a line alone does not canonically choose a spatial three-space. After choosing such a split, the nonzero vectors have the count $7+8$, but the coordinate split is additional structure.

The zero vector is an origin, not a one-dimensional time axis. A translation $x\mapsto x+t$ is period two and has no oriented sign over $\F_2$, since $-t=t$. It cannot replace the integer-valued oriented winding $\tau_{\rm Th}$ defined in Section~\ref{sec:time}. A physical $3+1$ correspondence would need a native selection rule, compatible history orientation, interval structure, dynamics and a controlled physical interpretation.

\subsection{Maxwell as a specified reduction target}
The physical comparison is the coupled electric/magnetic system with divergence constraints, curl evolution and charge continuity, not an assignment of the six letters WXYZTI to six Cartesian field components \citep{feynmanmaxwell}. In a candidate finite model, oriented incidence operators, source variables and an update law must first be defined. Their constraints must be preserved under update, and the limiting observables must reproduce the electromagnetic equations with a specified constitutive law and calibration.

A generic pair of coupled channels or a wave equation is insufficient to establish this reduction. Electric charge current is not identified with information current by terminology, and the $\mathcal K_{120}$ parity bit is not a derived electromagnetic phase. A successful native update should also satisfy the existing event-preserving and locality hypotheses before comparing its characteristic propagation to $c_{\rm Th}=1$.
